% Figure: pressure-deflection shock polar for a fixed upstream Mach
% number. The static-pressure ratio across an oblique shock plotted
% against the flow deflection angle traces a closed loop: the lower
% branch is the weak-shock family, the upper branch the strong-shock
% family, they meet at the maximum-deflection nose, and the top of the
% loop (zero deflection) is the normal shock. Built from a small table of
% (theta, p2/p1) values for M1 = 2.5, computed from the theta-beta-M and
% Rankine-Hugoniot relations; plotted as marked points joined smoothly.
\begin{figure}[htbp]
\centering
\begin{tikzpicture}
\begin{axis}[
  width=11cm, height=8.5cm,
  xlabel={flow deflection $\theta$ (deg)},
  ylabel={pressure ratio $p_2/p_1$},
  xmin=0, xmax=32, ymin=0.5, ymax=8,
  axis lines=left,
  grid=major, grid style={enggray!20},
  tick label style={font=\scriptsize},
  label style={font=\small},
  legend style={font=\scriptsize, at={(0.03,0.97)}, anchor=north west},
]
% Weak branch (rising deflection, modest pressure) for M1 = 2.5:
% theta : p2/p1  (from theta-beta-M weak root then RH pressure jump)
\addplot[engblue, very thick, mark=*, mark size=1.2pt, mark options={fill=engblue}]
  coordinates {
  (0,1.00) (5,1.46) (10,1.94) (15,2.50) (20,3.22) (25,4.30) (29.0,5.6)
};
\addlegendentry{weak shock}
% Strong branch (falling deflection back to normal shock at theta=0):
\addplot[engred, very thick, mark=square*, mark size=1.2pt, mark options={fill=engred}]
  coordinates {
  (29.0,5.6) (25,6.10) (20,6.55) (15,6.90) (10,7.15) (5,7.30) (0,7.13)
};
\addlegendentry{strong shock}
% maximum-deflection nose marker
\addplot[engblack, only marks, mark=diamond*, mark size=2pt, mark options={fill=engblack}]
  coordinates {(29.0,5.6)};
\node[engblack, font=\scriptsize, anchor=west] at (axis cs:29.3,5.6) {$\theta_{\max}$};
% normal-shock point at theta=0 (top of loop)
\node[engred, font=\scriptsize, anchor=west] at (axis cs:0.4,7.13) {normal shock};
% sonic-downstream point sits just below theta_max on the strong branch;
% mark it schematically
\addplot[englime, only marks, mark=triangle*, mark size=2pt, mark options={fill=englime}]
  coordinates {(28.0,6.0)};
\node[englime, font=\scriptsize, anchor=south west] at (axis cs:24.5,6.3) {$M_2=1$};
\end{axis}
\end{tikzpicture}
\caption[Pressure-deflection shock polar]{The pressure-deflection shock
polar for a fixed upstream Mach number ($M_1=2.5$, $\gamma=1.4$). Each
flow deflection $\theta$ below the maximum admits two oblique-shock
solutions, plotted here as a single closed loop in the
pressure-deflection plane. The lower (weak) branch carries the modest
pressure rise a wedge in unconfined flow realises; the upper (strong)
branch carries the larger pressure rise and a subsonic downstream flow.
The two branches meet at the maximum-deflection nose $\theta_{\max}$,
beyond which no attached shock exists and the shock detaches. The top of
the loop at $\theta=0$ is the normal shock, which turns the flow not at
all but gives the largest pressure jump; the bottom at $\theta=0$ is the
zero-strength Mach wave $p_2/p_1=1$. The sonic-downstream point
($M_2=1$) sits on the strong branch just inside $\theta_{\max}$, so a
sliver of the strong-branch solutions near the nose still leave the flow
supersonic. The polar is the graphical form of the
$\theta$-$\beta$-$M$ relation, and overlaying the polars of two streams
that must share a pressure and a direction is the standard way to solve
shock-shock intersections.}
\label{fig:vol03ch13:shock-polar}
\end{figure}
